% Shrnutí okruhu Dobývání znalostí — jedna otázka = jedna strana.
% Kondenzát z 03-dobyvani-znalosti.tex; slouží k opakování těsně před zkouškou.
\documentclass[10pt,a4paper]{article}
\usepackage[utf8]{inputenc}
\usepackage[T1]{fontenc}
\usepackage{lmodern}          % vektorové fonty — nutné pro expansion v microtype
\usepackage[czech]{babel}
\usepackage[margin=1.3cm,top=1.2cm,bottom=1.3cm]{geometry}
\usepackage{amsmath,amssymb}
\usepackage{parskip}
\usepackage{enumitem}
\setlist{nosep,leftmargin=1.1em,labelsep=0.35em,topsep=1pt}
\usepackage{multicol}
\setlength{\columnsep}{5.5mm}
\usepackage{booktabs}
\usepackage{xcolor}
\definecolor{linkblue}{RGB}{20,60,150}
\definecolor{defbg}{RGB}{240,244,252}
\definecolor{defframe}{RGB}{100,130,190}
\definecolor{headcol}{RGB}{25,70,140}
\definecolor{markcol}{RGB}{176,84,0}
\usepackage[most]{tcolorbox}
\usepackage{microtype}
\usepackage[colorlinks=true,linkcolor=linkblue,urlcolor=linkblue]{hyperref}
\usepackage{fancyhdr}

\pagestyle{fancy}
\fancyhf{}
\renewcommand{\headrulewidth}{0.4pt}
\fancyhead[L]{\footnotesize\textsc{Dobývání znalostí --- shrnutí ke SZZ}}
\fancyhead[R]{\footnotesize\thepage}

% Nadpis otázky: číslo + doslovné znění věty oficiálního okruhu.
\newtcolorbox{qbox}[1]{colback=defbg,colframe=defframe,boxrule=0.7pt,arc=1.2mm,
  left=2.5mm,right=2.5mm,top=1.2mm,bottom=1.2mm,#1}
\newcommand{\qpage}[2]{%
  \begin{qbox}{}
  \textbf{\large Otázka #1.}\quad\textbf{#2}
  \end{qbox}
  \vspace{-0.3em}}
% Vnitřní mezititulek.
\newcommand{\hh}[1]{\par\smallskip\noindent{\bfseries\color{headcol}#1}\par\nobreak\vspace{-0.2em}}
\newcommand{\tm}[1]{\textbf{#1}}
\newcommand{\en}[1]{\emph{#1}}

\begin{document}
\raggedcolumns   % nestahovat/neroztahovat sloupce svisle

%=====================================================================
% Strana 1 — rozcestník
%=====================================================================
\thispagestyle{empty}
\begin{center}
{\LARGE\bfseries Dobývání znalostí}\\[0.4em]
{\large Shrnutí okruhu ke státní závěrečné zkoušce --- jedna otázka na jednu stranu}\\[0.6em]
{\small SZZ Umělá inteligence, MFF UK \quad$\cdot$\quad srpen 2026}
\end{center}

\begin{qbox}{}
\small\raggedright
\textbf{Oficiální znění okruhu} (zdroj pravdy, 2025/2026):\\
{\footnotesize\url{https://www.mff.cuni.cz/cs/studenti/bc-a-mgr-studium/studijni-plany/2025-2026/informatika/mgr/umela-inteligence}}
\end{qbox}

\footnotesize
Sedm vět oficiálního znění okruhu, sedm stran tohoto textu. Podrobný výklad je
v~\path{03-dobyvani-znalosti.pdf}, anglicky v~\path{ENG/03-data-mining.pdf}.

\begin{enumerate}[leftmargin=2em,itemsep=1pt]
\item Základní paradigmata dobývání znalostí. \hfill \textit{str.~2}
\item Příprava dat; výběr atributů a~metody pro analýzu jejich relevance. \hfill \textit{str.~3}
\item Metody pro dobývání znalostí; asociační pravidla, přístupy založené na principu učení
  s~učitelem a~klastrová analýza. \hfill \textit{str.~4}
\item Metody pro extrakci charakteristických diskriminačních pravidel a~měření jejich
  zajímavosti. \hfill \textit{str.~5}
\item Reprezentace, vyhodnocování a~vizualizace získaných znalostí. \hfill \textit{str.~6}
\item Modely pro analýzu sociálních sítí; míry centrality, detekce komunit. \hfill \textit{str.~7}
\item Praktické využití technik pro dobývání znalostí a~analýzu sociálních sítí.
  \hfill \textit{str.~8}
\end{enumerate}

\begin{multicols}{2}
\footnotesize

\hh{Průřezová témata --- ptají se napříč otázkami}
\begin{itemize}
\item \tm{Přesnost vs.\ srozumitelnost}: stromy a~pravidla jsou čitelné, SVM a~sítě přesnější
  a~neinterpretovatelné; typ úlohy (klasifikace / popis / nugety) rozhoduje, které kritérium
  vyhraje.
\item \tm{Nevyvážená data}: accuracy je zavádějící $\Rightarrow$ precision/recall, $F_1$,
  ROC/AUC, lift; váhy chyb nebo různé pravděpodobnosti výběru už při přípravě dat.
\item \tm{Škály atributů} (nominální, ordinální, intervalová, poměrová) určují přípustné
  operace, metriky i~volbu metody --- to je nejčastější \uv{proč} u~doplňujících otázek.
\item \tm{Transformace vs.\ selekce atributů}: PCA (nové atributy, ztráta interpretace, nutné
  znát všechny původní) vs.\ filtry a~wrappery (podmnožina, interpretace zůstává).
\item \tm{Složitost}: AOI $O(na)$, indukce stromu $O(np\log p)$, PAM $O(kn^2d)$,
  Girvanové--Newman $O(m^2n)$, Louvain $O(n\log n)$; NP-těžké: hledání klik, vyvážený řez
  s~vahami, maximalizace vlivu (hladově 63~\% optima).
\item \tm{Proces, ne algoritmus}: CRISP-DM je cyklický, příprava dat spotřebuje 60--80~\% času,
  nasazení není konec (drift, soukromí, spravedlnost).
\end{itemize}

\columnbreak

\hh{Co se chce spočítat u~tabule}
\begin{itemize}
\item \tm{Support, confidence, lift} z~malé tabulky transakcí; jeden krok Apriori
  (spojení $+$ prořezání).
\item \tm{Entropie a~informační zisk} pro volbu atributu; příklad s~úvěrem (příjem 0,5747;
  konto 0,6667; pohlaví 0,9183; nezaměstnaný 0,825), Gini index.
\item \tm{Naivní Bayes} na malé tabulce (0,1042 vs.\ 0,0416).
\item \tm{$\chi^2$} na čtyřpolní tabulce ($42{,}857>3{,}84$), stupně volnosti $(R-1)(S-1)$.
\item \tm{Kvartily, IQR} a~mez odlehlosti $\pm 1{,}5\,\mathrm{IQR}$, box plot.
\item \tm{t-weight a~d-weight} z~tabulky generalizovaných záznamů.
\item \tm{Matice záměn}: $p$, $r$, $F_1$; konstrukce \tm{ROC} ze žebříčku a~úvaha o~AUC.
\item \tm{Centrality} na malém grafu (stupňová, mezilehlá, blízkostní, vlastní);
  \tm{klastrovací koeficient}; \tm{modularita} $Q$.
\item \tm{Test homofilie} (5 z~18 vs.\ $2pq=8$), test korelace (14~\% vs.\ 49~\%).
\item \tm{PageRank} jednou iterací; \tm{HITS} s~normalizací.
\end{itemize}

\hh{Nejčastější chyby}
\begin{itemize}
\item Tvrdit, že vysoká t-weight znamená vysokou d-weight (typičnost $\neq$ rozlišovací síla).
\item Zapomenout, že u~\tm{MS-Apriori} \emph{neplatí} uzavřenost dolů.
\item Míchat trénovací a~testovací data, resp.\ ladit parametry na testovací množině
  (na to je \emph{validační} množina).
\item Zaměnit \tm{bagging} (snižuje rozptyl, paralelní) a~\tm{boosting} (snižuje zkreslení,
  sekvenční, citlivý na šum).
\item U~PageRanku tvrdit, že potřebuje normalizaci (nepotřebuje --- na rozdíl od HITS).
\end{itemize}

\end{multicols}

\newpage

%=====================================================================
\qpage{1}{Základní paradigmata dobývání znalostí}
%=====================================================================
\begin{multicols}{2}
\footnotesize

\hh{Co je KDD}
\tm{Dobývání znalostí z~databází} (\en{knowledge discovery in databases}) $=$ netriviální proces
získávání \emph{implicitních}, dosud \emph{neznámých} a~\emph{potenciálně užitečných} informací
z~dat. Je \tm{interaktivní} (uživatel rozhoduje) a~\tm{iterativní} (fáze se opakují).
\emph{Data mining} je v~úzkém smyslu jen jedna fáze tohoto procesu, v~běžné mluvě celý proces.

\hh{Pět fází KDD}
\emph{výběr dat $\to$ předzpracování $\to$ transformace $\to$ dobývání (data mining) $\to$
interpretace a~vyhodnocení}. Příprava dat (první tři fáze) spotřebuje
\tm{60--80~\% času} celého projektu --- to je nejčastěji ověřovaný fakt.

\hh{Kořeny oboru}
Umělá inteligence (strojové učení), \tm{databáze} (velké objemy, dotazování) a~\tm{statistika}
(inference, testování), motivované potřebou \emph{podpory rozhodování}.
Odlišení od \tm{OLAP}: OLAP odpovídá na \emph{předem položené} otázky (agregace, kostky,
drill-down), KDD hledá vzory, \emph{o~nichž předem nevíme}.
Odlišení od klasické statistiky: KDD pracuje s~daty, která nebyla sbírána pro daný účel,
a~hledá hypotézy, nejen je testuje.

\hh{Tři typy úloh}
\begin{itemize}
\item \tm{klasifikace a~predikce} --- rozhoduje \emph{přesnost}, srozumitelnost je bonus;
\item \tm{popis} (deskripce) --- rozhoduje \emph{srozumitelnost} a~\emph{plné pokrytí} dat;
\item \tm{hledání nugetů} --- rozhoduje \emph{překvapivost}; nuget nemusí pokrývat téměř nic,
  stačí, že je zajímavý a~akceschopný.
\end{itemize}

\hh{Dělení úloh a~metod}
\tm{Deskriptivní} (popsat data: asociační pravidla, klastrování, sumarizace) vs.\
\tm{prediktivní} (odhadnout neznámou hodnotu: klasifikace, regrese).
\tm{S~učitelem} (známé cílové hodnoty) vs.\ \tm{bez učitele} (jen struktura dat) vs.\
\tm{poloučené} (málo označkovaných $+$ hodně neoznačkovaných dat).

\hh{Škály atributů}
\tm{Nominální} (jen rovnost --- barva), \tm{ordinální} (uspořádání --- známka),
\tm{intervalová} (rozdíly, ne poměry --- teplota ve~$^\circ$C),
\tm{poměrová} (i~poměry, absolutní nula --- příjem).
Škála určuje \emph{přípustné operace} (nelze počítat průměr nominálního atributu),
volbu metriky i~metody. Typické transformace: nominální atribut $\to$ $v$ binárních atributů;
poměrový atribut s~exponenciálním chováním $\to$ logaritmovat.

\hh{Metodologie}
\tm{SEMMA} (SAS): \emph{Sample, Explore, Modify, Model, Assess}.
\tm{CRISP-DM} (6 fází): porozumění problému $\to$ porozumění datům $\to$ příprava dat $\to$
modelování $\to$ vyhodnocení $\to$ nasazení. Pořadí \emph{není striktní} a~proces je
\emph{cyklický}. \tm{ASUM-DM} (IBM) $=$ CRISP-DM $+$ projektové řízení a~provoz.

\hh{Nejčastější doplňující otázka}
\emph{Proč pořadí fází CRISP-DM není striktní?} --- protože výsledky každé fáze určují další
kroky: zjištěná kvalita dat vede zpět k~přeformulování cíle, výsledek modelování zpět
k~jiné přípravě dat. Proces je smyčka s~návraty, ne lineární postup.

\end{multicols}
\newpage

%=====================================================================
\qpage{2}{Příprava dat; výběr atributů a~metody pro analýzu jejich relevance}
%=====================================================================
\begin{multicols}{2}
\footnotesize

\hh{Cíl přípravy dat}
Vybrat relevantní data a~reprezentovat je tak, aby je zvolený algoritmus zvládl; výsledkem je
\emph{jedna} tabulka \emph{objekty $\times$ atributy} (denormalizace, agregace přes vztahy 1:N).

\hh{Chybějící hodnoty}
Možnosti: \emph{ignorovat objekt} (ztráta dat, riziko zkreslení) / \emph{nová hodnota
\uv{nevím}} / \emph{doplnit existující hodnotu} (modus, poměrně podle rozdělení, libovolně) /
\emph{doplnit modelem} (regrese, $k$-NN, EM).
Umět říct, kdy je která volba vhodná: je-li \emph{absence sama informací} (nevyplněný příjem
u~žadatele o~úvěr), patří jí vlastní kategorie \uv{nevím}.

\hh{Odlehlé hodnoty}
$\mathrm{IQR}=Q_3-Q_1$; odlehlá hodnota leží mimo
$[\,Q_1-1{,}5\,\mathrm{IQR},\; Q_3+1{,}5\,\mathrm{IQR}\,]$.
Umět spočítat kvartily ze vzestupně seřazených dat (pozice $n/4$ a~$3n/4$, zaokrouhlení nahoru)
a~nakreslit \tm{box plot} (medián, kvartily, vousy, body mimo).
Odlehlá hodnota může být chyba měření (odstranit) i~nejzajímavější případ (detekce podvodů).

\hh{Standardizace}
Desítkové škálování, \tm{min-max} do $[0,1]$ nebo $[-1,1]$,
\tm{z-skóre} podle \emph{střední absolutní odchylky} (robustnější k~odlehlým hodnotám),
\tm{std-skóre} podle směrodatné odchylky.
Motivace: bez standardizace ovládne vzdálenost atribut s~největším rozsahem (příjem
v~korunách vs.\ počet dětí).

\hh{Diskretizace}
\tm{Ekvidistantní} (stejně široké intervaly --- citlivá na odlehlé hodnoty),
\tm{ekvifrekvenční} (stejně obsazené), \tm{s~využitím třídy} (dělení podle entropie nebo
$\chi^2$ --- hranice tam, kde se mění rozdělení tříd).
\tm{Fuzzy diskretizace}: rozostřené hranice, jedné hodnotě odpovídají dvě kategorie se součtem
charakteristických funkcí 1; váha fuzzy-objektu je součin $\prod_j \mu_j$.
Rizika diskretizace: ztráta informace a~vznik \emph{kontradikcí} (shodné vzory s~různou třídou).
\tm{KEX} seskupuje hodnoty kategoriálního atributu podle $\chi^2$ testu proti \emph{apriornímu}
rozdělení tříd (procedury \emph{Assign} a~\emph{Group}).

\hh{Redukce dat}
\emph{Příliš mnoho objektů} $\to$ vzorkování (pozor na \tm{reprezentativnost}), efektivnější
uložení, více modelů nad podmnožinami. \emph{Nevyvážené třídy} $\to$ váhy chyb nebo různé
pravděpodobnosti výběru (převzorkování minority, resp.\ podvzorkování majority).
\emph{Příliš mnoho atributů} $\to$ dvě různé cesty:
\begin{itemize}
\item \tm{transformace} (\tm{PCA}): méně \emph{nových} atributů jako lineární kombinace
  původních; nutné však znát \emph{všechny} původní atributy a~mizí interpretace;
\item \tm{selekce}: podmnožina \emph{původních} atributů, interpretace zůstává.
\end{itemize}

\hh{Filtry vs.\ wrappery}
\tm{Filtr} počítá charakteristiku pro každý atribut \emph{samostatně}, nezávisle na modelu ---
rychlý, ale nezachytí \emph{interakce} atributů.
\tm{Wrapper} postaví model nad podmnožinou atributů a~vyhodnotí jej --- drahý, zohledňuje
interakce a~konkrétní algoritmus.
Prohledávání podmnožin \emph{bottom-up} (přidávání atributů) nebo \emph{top-down} (odebírání).

\hh{Filtrační kritéria}
$\chi^2$ (větší $=$ lepší), entropie $H(A)$ (menší $=$ lepší), vzájemná informace
$\mathrm{MI}(A,C)$, informační míra závislosti $\mathrm{ID}(A,C)=\mathrm{MI}(A,C)/H(C)$
(větší $=$ lepší).
Umět zdůvodnit \emph{normalizaci} $\mathrm{ID}$: bez ní jsou zvýhodněny atributy s~mnoha
hodnotami (extrém: identifikátor rozdělí data na jednoprvkové skupiny).

\hh{$\chi^2$ test a~regrese}
\tm{$\chi^2$ test} nezávislosti: $(R-1)(S-1)$ stupňů volnosti, podmínka použitelnosti
$r_k s_l/n \ge 5$; pro malé četnosti \tm{Fisherův exaktní test}.
Umět projít výpočet na čtyřpolní tabulce (příklad s~úvěrem: $\chi^2 = 42{,}857 > 3{,}84$
$\Rightarrow$ zamítáme nezávislost).
\tm{Regrese a~korelace}: metoda nejmenších kvadrátů, $\beta$ jako kovariance dělená rozptylem,
$\alpha=\bar y-\beta\bar x$; vícerozměrně $\beta=(X^{\!\top}X)^{-1}X^{\!\top}Y$.
Korelační koeficient $\rho=\sigma_{xy}/(\sigma_x\sigma_y)$ měří \emph{jen lineární} závislost
--- nulová korelace neznamená nezávislost.

\end{multicols}
\newpage

%=====================================================================
\qpage{3}{Metody pro dobývání znalostí; asociační pravidla, přístupy založené na principu
učení s~učitelem a~klastrová analýza}
%=====================================================================
\begin{multicols}{2}
\footnotesize

\hh{Asociační pravidla --- míry}
\tm{Support} $=$ podíl transakcí obsahujících $i$ i~$j$; \tm{confidence} $=$ podíl transakcí
s~$i$ i~$j$ \emph{mezi transakcemi s~$i$}; \tm{lift} $=p(i\wedge j)/(p(i)p(j))$.
Lift $<1$ $\Rightarrow$ pravidlo je horší než náhoda a~negace závěru může být lepší pravidlo.
\emph{Umět spočítat všechny tři z~malé tabulky transakcí.}

\hh{Apriori a~varianty}
\tm{Apriori} využívá \tm{uzavřenost dolů} (každá podmnožina frekventované množiny je
frekventovaná) a~prohledává \emph{do šířky}. \textsc{Apriori-Gen} $=$ krok \emph{spojení}
(kombinace shodné v~$k-2$ kategoriích) $+$ krok \emph{prořezání} (chybí-li některá podkombinace
délky $k-1$). Prostor pravidel je $O(2^m)$, MBA typicky vyprodukuje \emph{nadměrné množství}
pravidel.
\tm{Taxonomie a~virtuální položky}: vzácné položky brát výše v~taxonomii, časté níže; virtuální
položky nesou informaci o~transakci (den, čas, zákazník) --- pozor na redundanci pravidel.
\tm{MS-Apriori} (vzácné položky): každá položka má vlastní $\mathrm{MIS}$, minsup pravidla je
\emph{minimum} MIS jeho položek, navíc omezení rozdílu supportů $\varphi$.
\emph{Uzavřenost dolů přestává platit} $\Rightarrow$ položky se setřídí vzestupně podle MIS,
místo $F_1$ se použije množina \emph{semen} $L$ a~\texttt{tailCount} při generování pravidel.
\emph{Umět uvést protipříklad k~uzavřenosti dolů.}
\tm{CAR}: transakce označkované třídou, pravidlo $X\to y$; dolují se \emph{přímo v~jednom
kroku} přes pravidlové položky $(\mathit{condset},y)$; různým třídám lze dát různý minsup
(101~\% $=$ třídu zakázat).
\tm{Sekvenční vzory}: sekvence $=$ uspořádaný seznam množin položek, \emph{velikost} $=$ počet
prvků, \emph{délka} $=$ počet položek; \tm{GSP} pracuje obdobně jako Apriori. Pozor:
$\langle\{3\}\{8\}\rangle$ a~$\langle\{3,8\}\rangle$ se navzájem \emph{neobsahují}.
\tm{FP-Growth}: 2 průchody daty, \emph{žádní kandidáti}; FP-strom sdílí \emph{prefixy} (pořadí
položek podle klesajícího supportu je heuristika) $+$ spojové seznamy položek; extrakce zdola
nahoru \emph{rozděl a~panuj} přes prefixové podstromy a~podmíněné FP-stromy. Nejhorší případ:
unikátní transakce $\Rightarrow$ strom větší než data.

\hh{Rozhodovací stromy}
\tm{TDIDT} (rozděl a~panuj): atribut se volí podle \emph{minimální vážené entropie}
$H(A)=\sum_v \frac{|A(v)|}{|D|}H(A(v))$, ekvivalentně maximálního \emph{informačního zisku}.
Složitost indukce $O(np\log p)$, klasifikace $O(q\log p)$. Strom lze přepsat na pravidla.
\tm{ID3} (Quinlan 1986) --- hladový, informační zisk.
\tm{C4.5} --- chybějící a~spojité hodnoty, prořezávání validací (\en{reduced-error pruning},
\en{subtree raising}), \en{rule post-pruning}, odhad přesnosti dolní mezí
$\mathit{acc}-1{,}96\sqrt{\mathit{acc}(1-\mathit{acc})/n}$, \emph{poměrný} informační zisk
(penalizuje nadměrné větvení). \tm{C5.0} --- komerční, boosting.
\tm{CART} --- \emph{binární} stromy, Gini $G(v_i)=1-\sum_j p_j^2$ (nižší $=$ lepší), kritérium
$\Phi=2P_LP_R\sum_j |P(C_j|t_L)-P(C_j|t_R)|$. \tm{CHAID} --- $\chi^2$, automatické seskupování
hodnot kategoriálních atributů.

\hh{Další metody učení s~učitelem}
\tm{Bayes}: hypotéza s~maximální aposteriorní pravděpodobností
$H_{MAP}=\arg\max_H P(E\mid H)P(H)$ (jmenovatel lze zanedbat); \emph{naivní}
předpoklad podmíněné nezávislosti $\Rightarrow \arg\max_{H_t}P(H_t)\prod_i P(E_i\mid H_t)$.
Zvládne nekompletní vzory; problém nulových pravděpodobností (Laplaceovo vyhlazení).
\tm{SVM}: maximální okraj $2/\|w\|$, primární úloha $\min \frac{\langle w\cdot w\rangle}{2}$
za podmínek $y_i(\langle w\cdot x_i\rangle+b)\ge 1$; Kuhn--Tuckerovy podmínky jsou u~konvexní
úlohy s~lineárními omezeními nutné \emph{i~postačující}; $\alpha_i>0$ jen pro \emph{podpůrné
vektory} (Wolfeho duál). \tm{Kernelový trik}: duál potřebuje jen skalární součiny
$\Rightarrow K(x,z)=\langle\phi(x)\cdot\phi(z)\rangle$ a~$\phi$ nikdy nevyčíslíme (Mercer). \emph{Měkký okraj}: pomocné proměnné $\xi_i\ge0$
v~podmínkách a~penalizace $C\sum_i\xi_i$ v~účelové funkci; $C$ řídí kompromis mezi šířkou
okraje a~chybami.
Omezení: dvě třídy, reálný prostor, neinterpretovatelné.
\tm{Logistická regrese}: $\ln\frac{p}{1-p}=\beta^{\!\top}x$, tedy $p=\sigma(z)$; učení
maximální věrohodností (křížová entropie), $\sigma'=\sigma(1-\sigma)\Rightarrow
\Delta\beta_j=\eta\sum_i (y_i-\sigma(z_i))x_{ij}$.
\tm{Diskriminační analýza}: $g_j(x)$ pro každou třídu, klasifikuje maximum; optimální
$g_j\propto P(C_j\mid x)$; normální rozdělení $\to$ kvadratická funkce, při \emph{shodných}
kovariančních maticích \emph{lineární} (stačí $\mu_j$ a~$S$).
\tm{ELM}: náhodné parametry skrytých neuronů, výstupní váhy \emph{jednou} pseudoinverzí
$\beta=H^{+}T$ (regularizace $I/C$); rychlé, bez lokálních minim.
\tm{Ensembly}: \tm{bagging} snižuje \emph{rozptyl} ($\sigma^2\to\sigma^2/k$ pro i.i.d.),
bootstrap obsahuje $\approx 63{,}2$~\% různých bodů ($1-1/e$), ideální komponenta je strom
(nízké zkreslení, vysoký rozptyl); zkreslení nesníží.
\tm{Náhodné lesy} řeší korelaci komponent: rozptyl $\rho\sigma^2+(1-\rho)\sigma^2/k$ ---
první člen na $k$ nezávisí; náhodnost přes bootstrap \emph{a} podmnožinu
$m\approx\sqrt{\#\text{příznaků}}$ v~každém uzlu.
\tm{AdaBoost} snižuje \emph{zkreslení}: váhy chybně klasifikovaných rostou,
$\alpha_t=\frac12\ln\frac{1-\epsilon_t}{\epsilon_t}$, běží \emph{sekvenčně}, citlivý na šum.

\hh{Klastrová analýza}
\tm{Metriky}: Hammingova/Manhattan ($L_1$), euklidovská ($L_2$), Čebyševova ($L_\infty$),
Minkowského ($L_q$ --- ostatní jsou její speciální případy), \tm{Mahalanobisova}
$d_M^2=(x_1-x_2)^{\!\top}S^{-1}(x_1-x_2)$ pro veličiny s~různým rozptylem. Před shlukováním
\emph{normalizovat}. Vzdálenost klastrů: nejbližší soused, nejvzdálenější soused, průměrná,
centroidní.
\tm{$k$-means} (Lloyd): náhodné rozdělení $\to$ centroidy $\to$ přesuny $\to$ opakuj;
cenu nikdy nezvýší, ale bez záruky aproximace $\Rightarrow$ \tm{$k$-means++} (pravděpodobnost
$\propto$ kvadrát vzdálenosti k~nejbližšímu centroidu), LocalSearch++.
\tm{$k$-medians} (Manhattan, medián po dimenzích, robustnější) vs.\ \tm{$k$-medoids}
(reprezentanti \emph{z~dat}, hill-climbing s~výměnami, $O(knd)$ na iteraci, libovolný typ dat):
\tm{PAM} testuje všech $k(n-k)$ výměn ($O(kn^2d)$), \tm{CLARA} vzorkuje (riziko špatného
vzorku), \tm{CLARANS} prohledává náhodně nad plnými daty (\emph{MaxAttempt}, \emph{MaxLocal}).
\tm{Hierarchické} zdola nahoru: každý vzor svůj klastr $\to$ slučuj dva nejbližší; výstup
\tm{dendrogram}, počet klastrů se volí až \emph{řezem}. \tm{CURE}: vzorek $\to$ $p$ částí
$\to$ nezávisle hierarchicky na $k'$ $\to$ hierarchicky na $k$ $\to$ přiřaď zbytek.
\tm{LVQ} $=$ shlukování \emph{s~učitelem} (vítězný vektor se přitáhne při souhlasu tříd,
odtlačí při nesouhlasu, $\mu$ klesá jako $\mu(k-1)/(k+1)$).
\tm{Mřížkové} (MAFIA): $p^d$ hyperkrychlí, prah hustoty $\tau$, klastry $=$ komponenty
souvislosti grafu buněk. \tm{DBSCAN}: \emph{jádrový} bod ($\ge\tau$ bodů v~radiu
$\varepsilon$), \emph{hraniční}, \emph{šumový}; klastry $=$ komponenty souvislosti jádrových
bodů. Zvládne klastry \emph{libovolného tvaru} a~sám detekuje odlehlé hodnoty.

\end{multicols}
\newpage

%=====================================================================
\qpage{4}{Metody pro extrakci charakteristických diskriminačních pravidel a~měření jejich
zajímavosti}
%=====================================================================
\begin{multicols}{2}
\footnotesize

\hh{Charakterizace vs.\ diskriminace}
\tm{Charakterizace} popisuje \emph{cílovou třídu} --- co je pro ni typické.
\tm{Diskriminace} ji \emph{porovnává} s~kontrastními třídami --- čím se od nich odlišuje.
Rozdíl je v~otázce: \emph{typičnost} $\neq$ \emph{rozlišovací síla}. Vlastnost může být pro
třídu typická (většina zákazníků jsou muži), a~přesto bezcenná pro rozlišení (muži tvoří
většinu i~v~ostatních třídách).

\hh{AOI --- indukce řízená atributy}
Postup \en{attribute-oriented induction}:
\begin{enumerate}
\item sběr relevantních dat dotazem (cílová, případně kontrastní třída);
\item přidání atributu \texttt{count} k~záznamům;
\item pro každý atribut \emph{odstranění} (má mnoho hodnot a~neexistuje k~němu koncepční
  hierarchie) nebo \emph{generalizace} (nahrazení nadřazenou hodnotou podle hierarchie);
\item agregace shodných záznamů (součet \texttt{count});
\item další generalizace, dokud relace není menší než prahová velikost $T$;
\item prezentace výsledné generalizované relace.
\end{enumerate}
Složitost $O(na)$ --- \emph{lineární} v~počtu záznamů i~atributů.
Prahy: \emph{attribute threshold} (kolik hodnot atributu se ještě snese) a~\emph{generalized
relation threshold} $T$.

\hh{Koncepční hierarchie}
\tm{Schématová} (dána strukturou dat: ulice $\to$ město $\to$ stát), \tm{seskupovací}
(explicitní výčet skupin hodnot), \tm{operační} (odvozená výpočtem: e-mail $\to$ doména),
\tm{pravidlová} (definovaná pravidly). Bez hierarchie lze atribut jen odstranit.

\hh{Dvě váhy}
$$\text{t-weight}=\frac{\mathrm{count}(q_a)}{\sum_i \mathrm{count}(q_i)}$$
Typičnost \emph{uvnitř} cílové třídy; součet t-weight přes všechny záznamy generalizované
relace je 100~\%. Slouží pro \tm{charakteristická} pravidla --- implikace \uv{$\Rightarrow$}
ve směru \emph{třída $\Rightarrow$ vlastnosti} (nutná podmínka).
$$\text{d-weight}=\frac{\mathrm{count}(q_a\in C_{\mathit{target}})}
{\sum_j \mathrm{count}(q_a\in C_j)}$$
Rozlišovací síla vůči kontrastním třídám (jmenovatel je součet přes \emph{všechny}
třídy včetně cílové). Slouží pro \tm{diskriminační} pravidla ---
\uv{$\Leftarrow$}, tedy \emph{vlastnosti $\Rightarrow$ třída} (postačující podmínka).
\tm{Popisné pravidlo} kombinuje obojí (\uv{$\Leftrightarrow$}) a~uvádí obě váhy.

\hh{Interpretace}
\emph{Umět spočítat obě váhy z~tabulky} a~vysvětlit, proč vysoká t-weight neznamená vysokou
d-weight: d-weight je nutné srovnávat s~\emph{apriorním} podílem cílové třídy. Je-li podíl
třídy v~datech 50~\% a~d-weight vlastnosti je také 50~\%, vlastnost nerozlišuje vůbec.

\hh{Zajímavost znalosti}
\tm{Objektivní} míry (počítá stroj): t-weight a~d-weight, support a~confidence, \tm{lift},
pokrytí, $\chi^2$, informační zisk, \emph{jednoduchost} (krátká pravidla).
\tm{Subjektivní} míry (rozhoduje uživatel): \tm{neočekávanost} (odporuje dosavadní představě)
a~\tm{akceschopnost} (\en{actionability} --- lze podle toho něco udělat).
Objektivní míry umí systém spočítat sám, ale teprve subjektivní rozhodují, zda je výsledek
\emph{znalost}, nebo notorieta --- stejné pravidlo může být pro jednoho uživatele objevem
a~pro druhého samozřejmostí.

\hh{Souvislosti}
Charakterizace a~diskriminace jsou \emph{deskriptivní} úlohy (otázka~1); vstupem je
tabulka po přípravě dat a~diskretizaci (otázka~2) plus koncepční
hierarchie, výstupem generalizovaná relace; zajímavost jako filtr
nadprodukce pravidel navazuje na asociační pravidla (otázka~3) a~na kritéria hodnocení
znalostí (otázka~5).

\end{multicols}
\newpage

%=====================================================================
\qpage{5}{Reprezentace, vyhodnocování a~vizualizace získaných znalostí}
%=====================================================================
\begin{multicols}{2}
\footnotesize

\hh{Formy reprezentace znalostí}
\tm{IF-THEN pravidla} --- nejsrozumitelnější, lokální (každé pravidlo lze číst samostatně);
\tm{rozhodovací strom} --- globální, čitelný do hloubky 4--5; \tm{rozhodovací tabulka};
\tm{matematická funkce} (SVM, sítě, regrese) --- kompaktní, ale neinterpretovatelná;
\tm{prototypy/centroidy} (klastrování); \tm{pravděpodobnostní model} (Bayes);
\tm{grafy a~sítě}; \tm{instance} ($k$-NN, CBR --- \uv{model} jsou samotná data).
Ke každé formě umět říct, z~jaké metody vzniká a~čím se za srozumitelnost platí.

\hh{Kritéria hodnocení metod}
Prediktivní přesnost, efektivita (čas konstrukce \emph{i} použití), robustnost (šum, chybějící
hodnoty), škálovatelnost, interpretovatelnost, kompaktnost modelu. Volba metody je vždy
kompromis mezi nimi --- viz typ úlohy z~otázky~1.

\hh{Odhad chyby}
\tm{Holdout} --- rozdělení na trénovací a~testovací množinu (velká data);
\emph{množiny se nesmí míchat}.
\tm{$k$-násobná křížová validace} --- data na $k$ částí, $k$ modelů; běžné je 5 a~10; přesnost
$=$ průměr $k$ přesností $+$ interval spolehlivosti s~$t_{N,k-1}$ (menší data).
\tm{Leave-one-out} --- $m$-násobná CV pro velmi malá data.
\tm{Validační množina} --- třetí část dat; na ní se ladí \emph{parametry} algoritmu, aby
testovací množina zůstala nezávislá.

\hh{Matice záměn a~míry}
$TP$, $FN$, $FP$, $TN$; přesnost $p=\frac{TP}{TP+FP}$, úplnost $r=\frac{TP}{TP+FN}$,
$F_1=\frac{2pr}{p+r}$ (harmonický průměr --- leží blíž k~\emph{menšímu} z~obou čísel, proto
netoleruje nulu ani v~jedné míře).
Umět vysvětlit, \emph{proč accuracy nestačí} u~nevyvážených dat: při 1~\% pozitivních případů
má triviální \uv{nedělat nic} 99~\% přesnost.

\hh{ROC a~AUC}
$\mathit{TPR}=\frac{TP}{TP+FN}$ (senzitivita $=$ recall pozitivní třídy) proti
$\mathit{FPR}=\frac{FP}{FP+TN}$; $\mathit{TNR}=$ specificita $=1-\mathit{FPR}$.
Křivka vede z~$(0,0)$ do $(1,1)$, diagonála $=$ náhodné hádání, levý horní roh $=$ ideál.
Konstrukce: vzory se seřadí podle skóre a~postupně se posouvá práh (dělení na každém případu).
\tm{AUC} $=1$ pro perfektní, $0{,}5$ pro náhodný klasifikátor; je potřebná proto, že se ROC
křivky dvou klasifikátorů mohou \emph{křížit} --- pak žádný není lepší ve všech prazích.

\hh{Skórování a~lift}
Místo třídy se přiřadí \emph{odhad pravděpodobnosti} (u~stromu confidence v~listu); vzory se
seřadí a~rozdělí do přihrádek (decilů). \tm{Lift křivka} porovnává zásah v~přihrádce
s~diagonálou náhodného výběru. Umět provést ekonomickou úvahu: kolik zásilek se vyplatí
poslat, když známe cenu zásilky a~výnos z~reakce.

\hh{Tři roviny vyhodnocování}
\begin{enumerate}
\item statistická kvalita modelu (přesnost, AUC) --- objektivní;
\item zajímavost jednotlivých vzorů (lift, $\chi^2$, jednoduchost) --- objektivní;
\item vyhodnocení uživatelem --- subjektivní.
\end{enumerate}
\tm{Čtyři kritéria zajímavosti znalosti}: \emph{srozumitelnost}, \emph{platnost},
\emph{užitečnost}, \emph{novost}.

\hh{Vizualizace}
Technika podle úlohy: histogram (rozdělení), box plot (odlehlé hodnoty), scatter a~SPLOM
(vztahy dvojic), paralelní souřadnice (mnoho dimenzí), heatmapa (matice), dendrogram
(hierarchie klastrů), projekce (PCA, t-SNE, UMAP), ROC a~lift (kvalita modelu).
\tm{Tři role v~KDD}: explorace dat (\tm{Anscombeho kvartet} --- shodné statistiky, různé grafy),
vizualizace modelu, prezentace výsledků.
Zásady (Tufte): nezkreslovat osami, poměr dat k~inkoustu, barva podle typu dat (kvalitativní
vs.\ sekvenční vs.\ divergující). Znát omezení: \emph{t-SNE a~UMAP zkreslují globální
vzdálenosti} --- vzdálenost mezi klastry v~projekci nelze interpretovat.

\end{multicols}
\newpage

%=====================================================================
\qpage{6}{Modely pro analýzu sociálních sítí; míry centrality, detekce komunit}
%=====================================================================
\begin{multicols}{2}
\footnotesize

\hh{Co je SNA}
Sociální vědy $+$ analýza sítí $+$ teorie grafů; zajímají ji \emph{vztahy}, ne entity a~jejich
atributy. Reprezentace: obrázek, seznam hran, matice sousednosti (u~orientovaných grafů
nemusí být symetrická). \tm{Ego síť} vs.\ \uv{celá} síť --- \tm{problém specifikace hranic}: žádná
studovaná síť není celá.

\hh{Síla vazeb}
Silné vs.\ slabé vazby; \tm{síla slabých vazeb} (Granovetter 1973 --- práci člověk najde přes
slabé vazby, protože nesou nové informace).
\tm{Homofilie} (podobní se spojují), \tm{tranzitivita} a~\tm{přemostění};
homofilie $+$ tranzitivita $\to$ \emph{kliky}.
\tm{Most} $=$ hrana, jejíž odstranění rozpojí koncové uzly; \tm{zkratkový most} $=$ vzdálenost
po odstranění vzroste nad 2. \tm{Překryv sousedství}
$O(i,j)=\frac{|N_i\cap N_j|}{|N_i\cup N_j|-2}$ --- čím větší, tím silnější vazba.

\hh{Čtyři míry centrality}
\begin{itemize}
\item \tm{stupňová} --- kolik lidí uzel \emph{přímo} dosáhne (lokální aktivita);
\item \tm{mezilehlá} --- na kolika nejkratších cestách leží; kde se síť \emph{rozpadne},
  kdo kontroluje tok;
\item \tm{blízkostní} --- převrácená průměrná délka nejkratších cest; rychlost \emph{šíření};
\item \tm{vlastní} --- $Ax=\lambda x$, jen \emph{největší} vlastní číslo dává nezáporný vektor;
  jak dobře je uzel spojen s~\emph{dobře propojenými}.
\end{itemize}
Umět interpretační tabulku a~vysvětlit, proč nestačí jedna míra (dva průměrné uzly mohou být
\emph{společně} důležitější než jeden s~nejvyšším stupněm).
\tm{Strukturní ekvivalence}: pozice lze vyměnit bez změny struktury sítě.

\hh{Koheze a~vlastnosti reálných sítí}
Reciprocita (jen orientované grafy), hustota $|E|/\binom{|V|}{2}$ (klika $=1$),
\tm{klastrovací koeficient} $C_i=\frac{2|\{e_{jk}\}|}{k_i(k_i-1)}$ a~jeho průměr,
\emph{průměr sítě} (\en{diameter}) $=$ nejdelší nejkratší cesta, průměrná vzdálenost,
excentricita. \tm{Malý svět} $=$ vysoké $C$ $+$ krátká průměrná cesta.
\tm{Šest vlastností reálných komplexních sítí}: efekt malého světa (Milgram, medián 6),
tranzitivita/klastrování, \emph{mocninné} rozdělení stupňů, odolnost (robustní k~náhodnému,
zranitelná cíleným odstraněním hubů), vzorce míchání (asortativní $=$ homofilie), komunitní
struktura.
\tm{Bezškálové sítě}: růst $+$ \tm{preferenční připojování} (rich-get-richer, Matoušův efekt);
\tm{centralizace} a~struktura \emph{jádro--periferie} (bow-tie).

\hh{Modely šíření vlivu}
Společné rysy: orientovaný graf, aktivní/neaktivní uzly, \emph{aktivovaný uzel nelze
deaktivovat}.
\tm{LTM} (prahový): $\theta_v$ uniformně z~$[0,1]$, aktivace při
$\sum_{\text{aktivní}} w_{u,v}\ge \theta_v\sum w_{u,v}$ --- pohled \emph{příjemce}.
\tm{ICM} (kaskádový): každý soused se aktivuje s~pravděpodobností $p_{v,u}$, ale
\emph{jen jednou} --- pohled \emph{vysílajícího}.
\tm{Maximalizace vlivu} je NP-těžká, hladový přístup dává $\ge 63$~\% optima.

\hh{Vliv, korelace, homofilie}
\tm{Test korelace}: podíl hran mezi různými hodnotami atributu vs.\ $2p(1-p)$ (kuřáci:
$14~\% < 49~\%$). Tři procesy, které korelaci vysvětlují: \emph{homofilie}, \emph{zmatení}
(společné prostředí), \emph{vliv}.
Logistický model $\ln\frac{p(a)}{1-p(a)}=\alpha\ln(a+1)+\beta$, $\alpha$ $=$ koeficient
sociální korelace; \tm{shuffle test} zpřehází \emph{časové značky} --- významná změna $\alpha$
je důkaz \emph{vlivu}.
\tm{Test homofilie}: heterogenních hran významně méně než $2pq$ $\Rightarrow$ homofilie,
významně více $\Rightarrow$ \emph{inverzní} homofilie (příklad: 5 z~18 vs.\ $2pq=8$ z~18).
\tm{Výběr vs.\ sociální vliv}: výběr $=$ charakteristiky řídí tvorbu vazeb, vliv $=$ vazby
formují \emph{mutabilní} charakteristiky (imutabilní: pohlaví, rasa, věk). Rozlišit lze jen
\emph{longitudinálně} (obezita: Christakis \& Fowler).
\tm{Afiliační sítě} (bipartitní osoby $\times$ fokusy) a~\tm{tři uzávěry}: tranzitivní
(společný přítel), \emph{fokální} (společný fokus $=$ výběr), \emph{členský} (přítel už ve
fokusu $=$ vliv); bázový model $T_{\text{base}}(k)=1-(1-p)^k$, empiricky klesající výnosy.
\tm{Schellingův model}: agenti dvou typů, spokojeni při $\ge t$ sousedech téhož typu
$\Rightarrow$ \emph{samovolná segregace} z~lokální homofilie.
\tm{Strukturní balance}: balancované trojúhelníky jsou $+{+}{+}$ a~$+{-}{-}$;
\emph{věta o~balanci} --- úplný balancovaný graf znamená všichni přátelé, nebo dvě skupiny
vnitřních přátel a~vzájemných nepřátel; \emph{slabá balance} (zakázáno jen $+{+}{-}$)
$\Rightarrow$ libovolný počet skupin.

\hh{Autority: HITS a~PageRank}
\tm{HITS} --- dotazově \emph{závislý}; kořenová množina $S$ $\to$ bázová $B$;
$a=L^{\!\top}h$, $h=La$, mocninná iterace s~\emph{nutnou normalizací};
$a_k=(L^{\!\top}L)^{k-1}a_0$ ($L^{\!\top}L$ $=$ ko-citační matice). Slabiny: spamovatelnost,
drift tématu, náklady v~době dotazu.
\tm{PageRank} --- dotazově \emph{nezávislý}, $P(i)=\sum_{(j,i)\in E}P(j)/O_j$,
\emph{normalizaci nepotřebuje} (celkový PageRank je konstantní). Problém \uv{pomalého úniku}
$\Rightarrow$ tlumicí faktor $d$ ($0{,}85$): $P(i)=(1-d)+d\sum P(j)/O_j$.
Markovský řetězec potřebuje matici \emph{stochastickou, ireducibilní a~aperiodickou} --- webový
graf nesplňuje ani jednu: stochasticitu kazí \emph{visící stránky} (nulový řádek --- řeší se
zvlášť: odstranit, nebo dát řádek $1/n$), ireducibilitu a~aperiodicitu spraví \emph{jediná}
úprava (odkaz odkudkoli kamkoli s~malou pravděpodobností).

\hh{Detekce komunit}
Definice komunity je \emph{subjektivní}; explicitní vs.\ implicitní skupiny.
\tm{Řezy}: minimální řez dává nevyvážené skupiny, s~vahami a~balancí je NP-těžký $\Rightarrow$
\emph{poměrový} a~\emph{normalizovaný} řez (dělení $|C_i|$, resp.\ $\mathrm{vol}(C_i)$).
\tm{Girvanové--Newman} --- divizivní, odstraňuje hrany s~nejvyšší \emph{mezilehlostí},
$O(m^2n)$; počet komunit podle maxima \tm{modularity}
$Q=\frac{1}{2m}\sum_{i,j}\bigl(A_{ij}-\frac{k_ik_j}{2m}\bigr)\delta(c_i,c_j)$,
$Q\in[-1,1]$, smysluplné komunity $Q\ge 0{,}3$.
\tm{Louvain} --- hladová aglomerativní maximalizace modularity ve dvou fázích (lokální přesuny
$\to$ supergraf), $O(n\log n)$; nevýhody: citlivost na pořadí, \emph{nesouvislé komunity},
lokální optima. \tm{Leiden} přidává fázi \emph{upřesnění rozdělení}.
\tm{Čtyři kategorie metod}: \emph{uzlově orientované} (kliky --- NP-těžké, \tm{CPM} pro
\emph{překrývající se} komunity, $k$-klika/klub/klan, $k$-jádro a~$k$-plex),
\emph{skupinově orientované} ($\gamma$-husté kvazikliky), \emph{síťově orientované}
(podobnost vrcholů $+$ $k$-means, MDS, blokové modely, spektrální shlukování přes laplacián
s~\emph{nejmenšími} vlastními čísly, maximalizace modularity s~\emph{největšími} --- vše jako
maticová faktorizace $X\approx S\Sigma S^{\!\top}$), \emph{hierarchicky orientované}
(divizivní GN, aglomerativní Louvain).
\tm{Tematické komunity} $K=(T,M)$ se mohou \emph{překrývat}; \emph{bipartitní jádra}
(vějíře $=$ huby, středy $=$ autority).

\hh{Predikce vazeb}
Obsahové míry (homofilie) vs.\ strukturní (tranzitivní uzávěr --- silnější).
\emph{Společní sousedé} $|S_i\cap S_j|$ (ignoruje stupně), \emph{Jaccard}
$\frac{|S_i\cap S_j|}{|S_i\cup S_j|}$, \emph{Adamic--Adar} $\sum_k \frac{1}{\log|S_k|}$
(váha klesá se stupněm společného souseda).
\tm{Katzova míra} $K_{ij}=\sum_t \beta^t W^{(t)}_{ij}$, maticově $K=(I-\beta A)^{-1}-I$
(pro dostatečně malé $\beta$) --- pro \emph{řídké} sítě
s~málo společnými sousedy.
\tm{Klasifikační pojetí}: pozitivně-neoznačkovaná úloha, vzorek negativ, logistická regrese;
klasifikátor se sám naučí relevanci příznaků a~zvládne \emph{asymetrické} příznaky
u~orientovaných sítí.

\end{multicols}
\newpage

%=====================================================================
\qpage{7}{Praktické využití technik pro dobývání znalostí a~analýzu sociálních sítí}
%=====================================================================
\begin{multicols}{2}
\footnotesize

\hh{Jak odpovídat}
Ke každé aplikaci umět říct \tm{trojici}: \emph{typ úlohy} (deskriptivní / prediktivní,
s~učitelem / bez učitele), \emph{použitá metoda}, \emph{způsob vyhodnocení}. Zkoušející
neověřuje výčet aplikací, ale schopnost přiřadit k~úloze metodu a~metriku.

\hh{Vlajkové aplikace dobývání znalostí}
\begin{itemize}
\item \tm{Analýza nákupního košíku} --- asociační pravidla; vyhodnocení supportem,
  confidence a~liftem; pozor na nadprodukci pravidel.
\item \tm{Detekce podvodů} (příklad zdravotní pojišťovny) --- stroj \emph{zaostřuje} pozornost
  experta, nenahrazuje jej; nevyvážená data, hodnotí se lift a~náklady chyb.
\item \tm{Kreditní skórování} --- rozhodovací stromy a~Bayes; klíčový požadavek
  \emph{srozumitelnosti} (rozhodnutí je nutné umět odůvodnit žadateli i~regulátorovi).
\item \tm{Segmentace zákazníků} --- klastrování; centroid $=$ typický zákazník; vyhodnocení
  interpretovatelností segmentů, ne přesností.
\item \tm{Churn a~cílený marketing} --- klasifikace na nevyvážených datech $\Rightarrow$
  ROC/AUC a~\emph{lift}; ekonomická úvaha o~počtu kontaktovaných zákazníků.
\item \tm{Klasifikace textu} --- SVM (vysoká dimenze, řídká data).
\item \tm{Predikce časových řad} --- posuvné okno (převod na tabulku objekty $\times$ atributy).
\item \tm{Web usage mining} --- sekvenční vzory nad logy.
\end{itemize}

\hh{Vlajkové aplikace SNA}
\begin{itemize}
\item \tm{Klíčoví hráči v~kriminálních a~teroristických sítích} --- míry centrality
  (mezilehlá pro spojky, vlastní pro vliv).
\item \tm{Doporučování přátel} --- predikce vazeb (společní sousedé, Adamic--Adar).
\item \tm{Optimalizace komunikačních sítí} --- koheze, úzká hrdla, mezilehlost.
\item \tm{Epidemiologie} (Rockdale County 1996) --- struktura kontaktní sítě vysvětlí šíření
  lépe než počet kontaktů.
\item \tm{Vakcinace cílená na huby} --- přímý důsledek vlastností bezškálových sítí.
\item \tm{Marketing a~kaskádové finanční krachy} --- modely šíření vlivu (LTM, ICM).
\item \tm{Segmentace zákazníků operátora} podle jazykových komunit --- Louvain nad sítí volání
  (2~mil.\ zákazníků; ukázka škálovatelnosti).
\end{itemize}

\hh{Klíčová vlastnost SF sítí pro praxi}
Odolnost proti \emph{náhodným} poruchám (síť vydrží až 80~\% odstraněných uzlů) vs.\
zranitelnost \emph{cíleným} útokem (stačí 5--15~\% hubů). Z~toho plyne jak strategie vakcinace
a~ochrany infrastruktury, tak nemožnost vymýtit internetové viry běžnými prostředky.

\hh{Nasazení není konec}
\tm{Monitoring driftu} (rozdělení dat se v~čase mění, model stárne), \emph{zpětná vazba modelu
do dat} (doporučení ovlivní chování, které se pak učí znovu), \tm{soukromí} (i~anonymizovaný
graf lze deanonymizovat porovnáním struktury), \tm{spravedlnost} (chráněný atribut lze
rekonstruovat z~ostatních, takže jeho vynechání nestačí), právní rámec ---
\tm{GDPR čl.~22} (automatizované rozhodování a~právo na vysvětlení) a~\tm{AI Act}.

\hh{Vazby na ostatní otázky}
Volba metody vychází z~typu úlohy (otázka~1), kvalita výsledku stojí na přípravě dat
(otázka~2), metriky a~ROC/lift pocházejí z~otázky~5, komunity a~centrality z~otázky~6.

\end{multicols}

\end{document}
